%TO CHECK!!

%Beyond technical contributions, this project supports broader social goals through the promotion of inclusive and accessible tools. 
%Open-source frameworks for power system analysis democratize advanced engineering capabilities, enabling participation from researchers 
%and engineers in diverse regions and institutions. For instance, the development of open and transparent electricity network 
%models demonstrates how access to modeling tools can empower under-resourced communities \cite{KirliPyPSA2021}. Open-source frameworks
%support the idea that knowledgment must be a social right and accessible for everybody.

%From a gender perspective, power systems engineering remains dominated by men in many contexts. Disseminating open methodologies 
%and encouraging diverse collaboration may help widen participation of underrepresented groups in STEM fields. 

%Finally, improving the reliability, stability, and efficiency of power systems has direct social benefits: fewer outages, better access to electricity, 
%and more resilient grids that support sustainable development goals.  


The development of power system analysis tools is not only a technical activity but also one with clear social implications,
particularly in the context of the ongoing energy transition. Modern electrical grids are critical infrastructures whose 
planning, operation, and regulation directly affect economic development, environmental sustainability, and social equity. 
As power systems evolve towards higher penetration levels of renewable and converter-based generation, the tools used to 
analyze stability and system behavior increasingly shape decision-making processes and, consequently, societal outcomes. 
For this reason, social and gender considerations are relevant even in technically focused research work.

The energy transition has been widely recognized as a multidimensional challenge that extends beyond technology deployment. 
According to the International Energy Agency (IEA), achie-\\ving long-term decarbonisation targets requires not only investment 
in renewable generation and grid infrastructure, but also a transformation of the workforce and the knowledge systems supporting 
energy planning and operation \cite{IEA2024Gender}. Engineering tools that enable secure and efficient integration of renewable 
energy sources contribute to system reliability and affordability, which are key elements of energy justice. 
In this sense, simulation frameworks such as the one developed in this work indirectly support broader societal goals 
by enabling informed technical decisions that reduce system risks and enhance the resilience of electricity networks.

Gender imbalance remains a structural issue in science, technology, engineering, and mathematics (STEM) fields, 
including power system engineering and software development. Multiple international studies show that women continue to be 
underrepresented in technical roles, particularly in areas related to electrical engineering, energy systems, and advanced 
software infrastructure \cite{UNESCOWomenSTEM, OECDGenderSTEM}. This imbalance limits diversity of perspectives in the 
design of technical tools and methodologies, potentially constraining innovation and reinforcing homogeneous problem-solving 
approaches. While this work does not directly address gender representation, it is developed within a sector where these 
structural challenges are present and should therefore be acknowledged.

Open-source software development constitutes a relevant mechanism for reducing some of the structural barriers associated 
with participation in technical fields. By providing unrestricted access to source code, documentation, and modelling 
methodologies, open-source platforms facilitate self-directed learning and enable contributions from a geographically and 
socially diverse community. This model lowers entry barriers related to licensing costs, institutional access, and proprietary 
knowledge, which disproportionately affect underrepresented groups. The European Commission has explicitly identified open 
science and open-source software as key enablers for inclusive innovation and equitable knowledge dissemination \cite{EUOpenScience}. 
In this context, the open-source nature of VeraGrid contributes to a more accessible technical ecosystem.

From a social perspective, transparent and openly available simulation tools also enhance accountability and reproducibility
in power system studies. Researchers, system operators, and students can independently verify assumptions, extend models, 
and adapt methodologies to local contexts. This openness supports capacity building in regions or institutions with limited 
access to commercial simulation tools, thereby contributing to a more equitable distribution of technical capabilities. 
Over time, such accessibility can support broader participation in energy system analysis and planning, which is particularly 
relevant in the context of rapidly evolving grids with high shares of renewable generation.

The social impact of advanced EMT simulation tools is also linked to their role in supporting reliable and secure electricity 
systems. Improved modelling of converter-dominated networks allows system operators and planners to anticipate stability issues, 
reduce the likelihood of large-scale disturbances, and design mitigation strategies. Reliable electricity supply is a 
prerequisite for social welfare, economic activity, and access to essential services. By contributing to the robustness of 
grid operation under high renewable penetration scenarios, this work indirectly supports societal resilience and the long-term 
sustainability of energy systems.

In a broader ethical context, engineers developing foundational simulation tools have a responsibility to consider the societal 
implications of their work. Promoting openness, transparency, and accessibility in technical development aligns with principles 
of responsible engineering and sustainable innovation. While open-source tools alone cannot resolve structural gender imbalances 
in STEM fields, they represent a necessary step towards more inclusive technical environments by reducing institutional and 
economic barriers. Future developments of the proposed framework could further enhance its social impact by incorporating 
educational resources, improving documentation accessibility, and encouraging contributions from a diverse user community.

Overall, this work contributes to the energy transition not only through its technical objectives but also by aligning with 
principles of inclusiveness, transparency, and social responsibility in engineering practice. By combining advanced simulation 
capabilities with an open-source development model, the proposed framework supports a more equitable and sustainable evolution 
of power system analysis tools.





